% Slide 1: OPSS Architecture and Execution Model
% Drop this frame into your Beamer deck.
% Requires: \usepackage{tikz}, \usetikzlibrary{arrows.meta, positioning, fit,
%            backgrounds, calc, decorations.pathreplacing}
% Requires color definitions (paste into preamble):
%   \definecolor{sensing}{HTML}{4A90D9}
%   \definecolor{vision}{HTML}{E67E22}
%   \definecolor{estimation}{HTML}{27AE60}
%   \definecolor{validation}{HTML}{8E44AD}
%   \definecolor{output}{HTML}{C0392B}
%   \definecolor{infra}{HTML}{7F8C8D}

\begin{frame}{OPSS Architecture and Execution Model}

% ── TOP: Pipeline dataflow ──────────────────────────────────────────
\begin{center}
\begin{tikzpicture}[
    >=Stealth,
    stage/.style={
      draw, rectangle, rounded corners=2pt,
      minimum width=1.6cm, minimum height=0.7cm,
      inner sep=2pt,
      align=center, font=\tiny\bfseries,
      text=white, line width=0.5pt},
    detail/.style={font=\fontsize{4}{5}\selectfont,
      text=black!75, align=center, inner sep=0pt},
    anno/.style={font=\fontsize{4}{5}\selectfont\itshape,
      text=gray!60!black},
    arr/.style={->, semithick, color=black!60},
  ]

  % --- Pipeline stages ---
  \node[stage, fill=sensing]                          (cam) {Sensing};
  \node[stage, fill=vision,     right=0.55cm of cam]  (vis) {Vision};
  \node[stage, fill=estimation, right=0.55cm of vis]  (est) {State Est.};
  \node[stage, fill=validation, right=0.55cm of est]  (val) {Valid.\,+\,Fusion};
  \node[stage, fill=output,     right=0.55cm of val]  (out) {Output};

  % --- Detail labels ---
  \node[detail, below=2pt of cam]  {RealSense D435\\1280$\times$720\\30\,fps RGB-D};
  \node[detail, below=2pt of vis]  {YOLOv8n FP16\\640$\times$360\\GPU};
  \node[detail, below=2pt of est]  {Kalman\,/\,UKF-NN\\6D state (m)\\NN residual};
  \node[detail, below=2pt of val]  {Physics bounds\\B$2_3$ weighted\\adaptive};
  \node[detail, below=2pt of out]  {Unix socket\\FastAPI dash.\\MyCobot\,280};

  % --- Arrows ---
  \draw[arr] (cam) -- node[above, anno] {frames} (vis);
  \draw[arr] (vis) -- node[above, anno] {dets}   (est);
  \draw[arr] (est) -- node[above, anno] {states} (val);
  \draw[arr] (val) -- node[above, anno] {fused}  (out);

\end{tikzpicture}
\end{center}

\vspace{-4pt}

% ── MIDDLE: Two-column layout ───────────────────────────────────────
\begin{columns}[T]

% --- Left column: Execution model ---
\begin{column}{0.52\textwidth}
\textbf{\small Execution Model}
\begin{itemize}[leftmargin=1.2em, itemsep=1pt, topsep=2pt]
  \item[\tiny$\blacktriangleright$] \textbf{2 daemon threads} --- camera capture +
        pipeline processing; web server on async event loop
  \item[\tiny$\blacktriangleright$] \textbf{Camera-paced} --- no internal timer;
        pipeline blocks on hardware vsync at 30\,fps
  \item[\tiny$\blacktriangleright$] \textbf{Bounded queue (depth\,=\,2)} ---
        \texttt{put\_nowait}; frames dropped under overload,
        not buffered
  \item[\tiny$\blacktriangleright$] \textbf{Sequential stages} --- all processing
        in one thread per frame; no intra-frame parallelism
  \item[\tiny$\blacktriangleright$] \textbf{Swappable tracker} --- Kalman $\leftrightarrow$
        UKF-NN behind duck-typed interface
        (\texttt{update(dets,\,t)\,$\to$\,List[ObjectState]})
\end{itemize}
\end{column}

% --- Right column: Timing budget ---
\begin{column}{0.46\textwidth}
\textbf{\small Timing Budget (est., 1 track)}

\vspace{2pt}
\begin{tikzpicture}[
    >=Stealth,
    bar/.style={minimum height=0.35cm, anchor=west,
      font=\fontsize{4.5}{5.5}\selectfont\bfseries, text=white,
      inner sep=2pt, rounded corners=1pt},
    lbl/.style={font=\fontsize{4.5}{5.5}\selectfont, text=black!70,
      anchor=west},
  ]

  % Scale: 1cm = 6ms.  Budget = 33.3ms = 5.55cm
  \def\W{5.55}

  % Background (full budget)
  \fill[black!6, rounded corners=1pt] (0,0) rectangle (\W, 0.35);

  % Detection: 18ms = 3.0cm
  \node[bar, fill=vision, minimum width=3.0cm] at (0,0) {YOLO 18\,ms};
  % Resize: 0.4ms = 0.067cm (too small to label)
  \fill[sensing!70, rounded corners=1pt] (3.0,0) rectangle (3.07, 0.35);
  % Track: 1ms = 0.167cm
  \fill[estimation!80, rounded corners=1pt] (3.07,0) rectangle (3.24, 0.35);
  % Post: 1.2ms = 0.20cm
  \fill[validation!70, rounded corners=1pt] (3.24,0) rectangle (3.44, 0.35);

  % Headroom brace
  \draw[<->, thick, color=black!45]
    (3.44, -0.12) -- (\W, -0.12)
    node[midway, below=1pt, lbl] {12.8\,ms headroom (38\%)};

  % 30fps deadline
  \draw[thick, red!55] (\W, -0.05) -- (\W, 0.45)
    node[above, font=\fontsize{4.5}{5.5}\selectfont\bfseries,
      text=red!55] {33.3\,ms};

  % Stage labels below
  \node[lbl] at (3.07, -0.38) {Track 1\,ms};
  \node[lbl] at (3.55, -0.38) {V+F+Bcast 1.2\,ms};

\end{tikzpicture}

\vspace{6pt}
\textbf{\small End-to-End Latency}

\vspace{1pt}
{\fontsize{6}{7.5}\selectfont
\begin{tabular}{@{}lr@{}}
  Sensor exposure + readout & $\sim$33\,ms \\
  Queue buffering (0--1 frame) & 0--33\,ms \\
  Pipeline processing & $\sim$20\,ms \\
  Socket tx & $<$0.1\,ms \\
  \hline
  \textbf{Glass-to-output} & \textbf{53--86\,ms} \\
\end{tabular}
}

\vspace{4pt}
{\fontsize{5.5}{7}\selectfont
\textit{Platform:} Jetson Orin Nano\\
6-core ARM A78AE\;\textbar\;1024 CUDA cores\;\textbar\;8\,GB unified}

\end{column}
\end{columns}

\end{frame}

% =====================================================================
% SPEAKER NOTES — key points to defend under faculty questioning
% =====================================================================
%
% Q: "Why sequential stages instead of a pipelined design?"
% A: Detection is a single indivisible GPU call (~18ms) that dominates
%    the frame budget. Pipelining would require inter-stage queues and
%    synchronization but the dominant cost cannot be overlapped with
%    itself. With 38% headroom at 1 track, the complexity is not
%    justified. (core.py:244-340)
%
% Q: "What happens under overload?"
% A: Frames are dropped at the queue boundary (camera.py:153-156,
%    put_nowait). The pipeline always processes the *most recent*
%    available frame. Worst case: effective FPS drops, tracker
%    covariance grows during longer inter-frame intervals, correctly
%    reflecting increased uncertainty. No latency accumulation.
%
% Q: "Why only 2 threads instead of a thread-per-stage?"
% A: Decoupling capture from processing prevents USB buffer corruption
%    under scheduling delays. Beyond that, adding threads adds sync
%    overhead without benefit: post-detection stages total ~2ms
%    combined. The web server runs on a separate async event loop
%    (Uvicorn) but is not in the control-critical path.
%
% Q: "Are the timing numbers measured?"
% A: These are analytical estimates from algorithmic complexity and
%    published NVIDIA TensorRT benchmarks for YOLOv8n FP16 on Orin Nano.
%    No empirical on-device measurements yet. Per-stage instrumentation
%    is planned (see Risk R3 / Section 7 of the PDR document).
%
% Q: "Why is the NN on CPU instead of GPU?"
% A: The NN has 611 parameters (15->32->3, two layers). At this scale,
%    GPU kernel launch latency (~10us) plus memory transfer overhead
%    exceeds the computation itself (<0.1ms as NumPy matmul). Keeping
%    YOLO as the sole GPU workload also eliminates kernel contention.
%    Additionally, PyTorch is unusable on this Jetson due to a
%    libcudnn version mismatch, so inference is pure NumPy.
%
% Q: "What's the worst-case latency?"
% A: 86ms (2.6 frames) — sensor readout (33ms) + 1 frame of queue
%    buffering (33ms) + processing (20ms). Irreducible floor is 53ms
%    (1.6 frames) when the queue is empty.
%
% Q: "How do you swap trackers?"
% A: Both MultiObjectKalmanFilter and MultiObjectUKFNN expose the same
%    duck-typed interface: update(detections, timestamp) -> List[ObjectState].
%    A single config flag (tracker_type = "kalman" | "ukf_nn") selects
%    at init time. Fallback: if UKF-NN import fails (torch unavailable),
%    pipeline auto-degrades to Kalman with a logged warning. (core.py:156-167)
%
% Q: "What are your external interfaces?"
% A: Three: (1) RealSense D435 via USB3 at 30Hz RGB-D (input),
%    (2) MyCobot 280 via Unix SOCK_DGRAM at /tmp/opss_cobot.sock,
%    30Hz JSON (output, unidirectional, non-blocking),
%    (3) Browser via FastAPI — REST, MJPEG, WebSocket (monitoring only,
%    not control-critical). Cobot/web failures are fault-isolated:
%    neither stalls the pipeline.
